% PathwayAnnotator 1.0.0: a report for Zenodo.
% Build: pdflatex, bibtex, pdflatex, pdflatex (main file pathwayannotator-1.0.0.tex).
\documentclass[10pt]{article}
\usepackage[a4paper,margin=2cm]{geometry}
\usepackage[T1]{fontenc}
\usepackage{lmodern}
\usepackage{microtype}
\usepackage{booktabs}
\usepackage{url}
\usepackage[hidelinks]{hyperref}
\usepackage[numbers,sort&compress]{natbib}
\setlength{\bibsep}{0pt plus 0.3ex}
\setlength{\parskip}{0.35em}
\setlength{\parindent}{0pt}

\title{\vspace{-1.5em}PathwayAnnotator: A Web-based SBML Annotation Tool for Biochemical Networks}
% Inside \author there can be no blank lines: use \\ for line breaks.
\author{Herbert M. Sauro$^{1,2}$\\[1ex]
  \small $^{1}$Department of Bioengineering, University of Washington, Seattle, WA, 98195-5061, USA\\
  \small $^{2}$eScience Institute, University of Washington, Seattle, WA, 98195-5061, USA
}
\date{\small Version 1.0.0, November 2026}

\begin{document}
\maketitle
\vspace{-1em}

\begin{abstract}
\noindent
PathwayAnnotator is a web application for annotating SBML and Antimony models. For metabolic models the tool will
suggest annotations for each species, reaction and compartment and shows the evidence for each suggestion. Other kinds of
model that include proteins or genes can also be annotated manually. The model
as a whole can also be annotated with its modelling framework, references, organism, creators and notes.
Accepted annotations are written into the model as MIRIAM and SBO terms. The application is written in Rust,
with a user interface built with the Leptos framework, and is compiled to WebAssembly so that it runs entirely
in the browser; suggestions come from a
knowledge base, built from BiGG Models, ChEBI and the curated models of BioModels, that is downloaded with
the page. On 133 held-out curated BioModels models with little overlap with the models used to build the
knowledge base, the first suggestion matched the curators' annotation for 90.8\% of species and 82.7\% of
reactions that could be graded. Suggestions the tool labels high confidence were correct in 99.3\% (species)
and 100\% (reactions) of cases, although these thresholds were chosen on the same test set.

\noindent\textbf{Availability:} \url{https://sys-bio.github.io/PathwayAnnotator/} (application) and
\url{https://github.com/sys-bio/PathwayAnnotator} (source, MIT or Apache-2.0).
\end{abstract}

\section{Background}
Annotations link the elements of a model to entries in external resources, so that a species named
\texttt{G6P} can be recognised as D-glucose 6-phosphate (ChEBI CHEBI:4170) by people and software. In SBML
\cite{hucka2003sbml} they are written as MIRIAM qualifier--resource pairs \cite{lenovere2005miriam} and
Systems Biology Ontology (SBO) terms \cite{courtot2011sbo}. Adding them by hand requires looking up every
element, and many published models are partly annotated or not annotated at all. Earlier tools such as
semanticSBML \cite{krause2010semanticsbml} and AMAS \cite{shin2023amas} address the same task. PathwayAnnotator
was written to combine automated suggestion with review by the modeller in a tool that needs no
installation and sends no model data to a server.

\section{Design}
\textbf{Architecture.} The application is written in Rust and compiled to WebAssembly; it is a static web page
and has no server component. The user interface uses Leptos (\url{https://leptos.dev}), a Rust framework
for web interfaces, in its client-side mode. Models are read and written with SBML (Levels 2 and 3) and Antimony
\cite{smith2009antimony} libraries written for the same project. The model never leaves the user's computer.

\textbf{Knowledge base.} A build program combines BiGG Models \cite{king2016bigg} (2,441 metabolites and
7,322 reactions, with their cross-references), ChEBI \cite{hastings2016chebi} release 257 (names, synonyms,
charge forms, and \emph{is\_a} relations used for broader and narrower matches), and the names that 1,075
curated BioModels models \cite{malik2020biomodels} use for annotated elements (a name is kept only when the
models using it agree on what it denotes). The result is a compressed file of about 2.2~MB. Display labels
are derived from the source names, preferring the common name of a molecule (3-phospho-D-glycerate) to the
name of its predominant charged form (3-phosphonato-D-glycerate). The data the published knowledge base is
built from is kept in the source repository, so that a given version can be rebuilt exactly.

\textbf{Suggestion engine.} Each candidate carries the evidence that produced it. Three sources of evidence
are combined: (i) exact lookup of element names and identifiers in the name index; (ii) a decoder for
names that are not found directly, which removes SBML prefixes, compartment words (\texttt{ATP\_mit},
\texttt{Glucose in}) and charges, and recognises database identifiers embedded in SBML identifiers; and
(iii) reaction context: a species is supported by knowledge-base reactions whose participants match the
reactions it takes part in in the model, and a reaction is supported by knowledge-base reactions whose
participants match its own. The first suggestion for each element is assigned one of three confidence
levels (high, likely, uncertain) from the kind of evidence, its score, and its lead over the next
candidate. Elements the user confirms or rejects
are fed back to the engine, which re-ranks the remaining suggestions; a reaction whose participants have all
been confirmed and which matches a single knowledge-base reaction becomes high confidence.

\textbf{Review and editing.} For each element the interface shows the suggestion, its confidence level,
synonyms, links to the database entries, and the evidence. Accepting writes a \texttt{bqbiol:is} term with
identifiers.org URIs and, where none is present, an SBO term. Users can change qualifiers, add or remove
terms, set SBO terms, and rename SBML identifiers, with all references (including kinetic laws, rules and
events) updated. Model-level annotation covers the modelling framework (SBO), references
(\texttt{bqmodel:isDescribedBy}), organism, creators, dates and notes. Every change can be undone.

\textbf{Optional online searches.} For elements the knowledge base does not cover, the user can search
ChEBI and the Gene Ontology \cite{go2023} through EBI's Ontology Lookup Service, reviewed UniProt entries
\cite{uniprot2023}, and Rhea \cite{bansal2022rhea}. Rhea is queried by the ChEBI identifiers of a
reaction's identified participants, or by enzyme name or EC number, and the results are ranked by how well
their participants match the model reaction. These searches are made only on request; the text sent is
shown before sending.

\section{Evaluation}
Every fifth curated BioModels model (215 of 1,075, chosen by model number) was held out, and a knowledge
base was built without them. The engine was given the names, identifiers, compartments and reaction
participants of each test model, but not its annotations. The curators' \texttt{bqbiol:is} and
\texttt{isVersionOf} annotations served as the answer key. Only elements annotated with an entry present in
the knowledge base can be graded: 1,267 of 3,569 annotated species and 641 of 4,095 annotated reactions. The
remainder are mostly proteins, complexes and pathway-level terms outside the scope of a metabolism
knowledge base. Some test models are close variants of training models; we therefore also report the 133
test models whose species names overlap by less than 50\% with every training model (610 species and 404
reactions graded). Table~\ref{tab:results} gives the contribution of each source of evidence on these
models.

\begin{table}[h]
\centering\small
\begin{tabular}{lcc}
\toprule
Evidence used (133 novel test models) & Species right first & Reactions right first \\
\midrule
Name lookup only & 85.1\% & \phantom{0}6.9\% \\
+ name decoder & 88.7\% & 32.7\% \\
+ reaction context (version 1.0.0) & 90.8\% & 82.7\% \\
\bottomrule
\end{tabular}
\caption{Share of graded elements whose first suggestion matched the curators' annotation. Coverage (an
element received any suggestion) was 96.7\% for species and 90.8\% for reactions; in the top three
suggestions, 92.6\% and 84.4\%.}
\label{tab:results}
\end{table}
\vspace{-0.5em}

Reaction names rarely identify a reaction across models, so name lookup alone recovered few reactions
on unseen models; most of the gain comes from matching participants. The confidence levels separate
reliable from unreliable suggestions: on the 133 models, high-confidence suggestions (70.7\% of graded
species with a suggestion, 30.2\% of reactions) were correct in 99.3\% and 100\% of cases, likely ones in
88.9\% and 89.6\%, and uncertain ones in 41.4\% and 46.7\%. When all species were given as confirmed, the
first reaction suggestion was correct for 84.9\% of reactions (82.7\% without); confirming reactions changed
species results little (91.0\% versus 90.8\%). On all 215 test models the first suggestion was correct for
91.9\% of species and 82.1\% of reactions.

The test models were held out of the knowledge base, but not out of development: the engine's rules, and
the round-number thresholds of the confidence levels, were chosen while inspecting results on this test
set. The figures above are therefore likely to be somewhat optimistic for models unlike those in BioModels.
An independent check needs annotated models from outside BioModels, or a three-way split of the curated
models.

\section{Limitations}
The automatic suggestions of version 1.0.0 cover small molecules and metabolic reactions. Proteins,
complexes, genes and mRNAs, and signalling and gene-regulation reactions (binding, phosphorylation,
transcription, degradation), are not suggested from the knowledge base; adding them is planned for version 2.
Models of signalling and gene regulation can nevertheless be annotated manually with the same tool: proteins
can be found with the UniProt search, terms and qualifiers added and edited, SBO terms set, and the model as
a whole annotated as for any other model. The knowledge base reflects
BiGG's coverage, which is weighted towards microbial and human central metabolism. The answer key consists
of curators' annotations, which are themselves sometimes broad (an EC number for a specific reaction) or
inconsistent between models, and elements annotated with entries outside the knowledge base are not graded,
so the figures describe performance on the annotatable part of the models. Rankings may change when the
underlying databases are updated. BiGG data from its legacy release is licensed for educational, research
and non-profit use; its notice is distributed with the knowledge base and shown in the application.

\section{Availability}
The application is available at \url{https://sys-bio.github.io/PathwayAnnotator/} and the source at
\url{https://github.com/sys-bio/PathwayAnnotator} under the MIT or Apache-2.0 licence. The repository
includes the benchmark program and its full results, and the data snapshot from which the published
knowledge base is built. Data sources: BiGG Models (CC BY 4.0 and the BiGG Models licence), ChEBI (CC BY 4.0),
BioModels (CC0).

\section*{Acknowledgements}
This work was supported by NIH Biomedical Imaging and Bioengineering award P41 EB023912 through HMS at the
Center for Reproducible Biomedical Modeling (\url{https://reproduciblebiomodels.org}). The software was
developed with the assistance of Claude Code (Anthropic).

{\small
\bibliographystyle{unsrtnat}
\bibliography{references}
}

\end{document}
